\section{First-Order Optimality Conditions}

In this section, we derive the first-order necessary optimality conditions associated with the 
stochastic optimal control problem introduced in Section~5.

\subsection{Variation with Respect to the Control}

Let \(u^\ast\in U_{ad}\) be an optimal control and let \(v\) be an admissible perturbation. 
Differentiating the Lagrangian \eqref{eq:lagrangian} with respect to the control variable yields

\begin{equation}
D_u \mathcal L(y,u,p)(v)
=
\E
\left[
\int_0^T
\left(
\alpha u - p,
v
\right)_H
dt
\right].
\label{eq:variation_control}
\end{equation}

A necessary condition for optimality is

\[
D_u \mathcal L(y^\ast,u^\ast,p^\ast)(v)=0,
\qquad
\forall v\in U_{ad}.
\]

Using \eqref{eq:variation_control}, we obtain

\begin{equation}
\E
\left[
\int_0^T
\left(
\alpha u^\ast-p^\ast,
v
\right)_H
dt
\right]
=
0,
\qquad
\forall v\in U_{ad}.
\label{eq:optimality_variational}
\end{equation}

Since \(v\) is arbitrary, it follows that

\begin{equation}
\alpha u^\ast-p^\ast=0.
\label{eq:control_relation}
\end{equation}

Consequently, the optimal control is given by

\begin{equation}
u^\ast
=
\frac{1}{\alpha}
p^\ast.
\label{eq:optimal_control}
\end{equation}

\subsection{Characterization of the Optimal Control}

Equation \eqref{eq:optimal_control} shows that the optimal control is directly proportional to 
the adjoint state. Therefore, once the adjoint variable has been computed, the optimal control 
can be obtained explicitly.

The parameter \(\alpha\) plays the role of a regularization coefficient. Large values of \(\alpha\) 
penalize the control effort and lead to smaller control amplitudes, whereas small values of 
\(\alpha\) allow stronger control actions.

\subsection{Optimality System}

Combining the state equation, the adjoint equation, and the optimality relation yields the complete 
optimality system.

The state equation is

\begin{equation}
dy
+
\left(
\mathbf b\cdot\nabla y
-
\kappa \Delta y
\right)dt
=
u\,dt
+
\sigma\,dW(t).
\label{eq:state_system}
\end{equation}

The adjoint equation is

\begin{equation}
-\frac{\partial p}{\partial t}
-\kappa \Delta p
-\nabla\cdot(\mathbf b,p)
=
y-y_d.
\label{eq:adjoint_system}
\end{equation}

The optimal control satisfies

\begin{equation}
u
=
\frac{1}{\alpha}p.
\label{eq:control_system}
\end{equation}

Together with the boundary and initial/final conditions,

\begin{equation}
\left\{
\begin{aligned}
y &=0
&&\text{on }
\partial D\times(0,T),
\
p &=0
&&\text{on }
\partial D\times(0,T),
\
y(\cdot,0)&=y_0,
\
p(\cdot,T)&=0,
\end{aligned}
\right.
\label{eq:boundary_conditions_system}
\end{equation}

Equations
\eqref{eq:state_system},
\eqref{eq:adjoint_system},
and
\eqref{eq:control_system}
form a coupled forward--backward system.

\subsection{Gradient Formula}

The adjoint state also provides a convenient expression for the gradient of the cost functional. 
Indeed,

\begin{equation}
\nabla J(u)
=
\alpha u-p.
\label{eq:gradient}
\end{equation}

This expression is particularly useful for gradient-based optimization algorithms.

For example, a gradient descent iteration reads

\begin{equation}
u^{k+1}
=
u^k
-
\rho
\left(
\alpha u^k-p^k
\right),
\label{eq:gradient_descent}
\end{equation}

where \(\rho>0\) denotes the step size.

The optimality system derived above constitutes the basis for the numerical approximation 
developed in the next section.
